%----------------------------------------------------------------------
\section{Detailed Verifications}\label{app:verifications}
%----------------------------------------------------------------------

This appendix collects component-by-component verifications
that support the main text's structural arguments. The main
text identifies \emph{why} each result holds; this appendix
verifies \emph{that} each component checks out.

%----------------------------------------------------------------------
\subsection{Triadic adjunction: full verification}
\label{app:adjunction-full}
%----------------------------------------------------------------------

We verify the $\sigma$-axis adjunction
$\kappa \leftarrow \sigma \to \tau$ in full detail.
The main text (Theorem \ref{thm:adjunction}) gives
the structural argument; here we check each component.

\emph{Functors on morphisms.}
The left adjoint $F: \mathcal{C}(S,\kappa') \to
\mathcal{C}(S,\kappa)$ acts on objects by inclusion
$\iota$ and on morphisms by
$F(d) = \iota(d)$. Composition preservation:
$F(d_2 \wedge d_1) = \iota(d_2 \wedge d_1) =
\iota(d_2) \wedge \iota(d_1) = F(d_2) \wedge F(d_1)$,
since $\iota$ is a ring homomorphism.

The right adjoint $U: \mathcal{C}(S,\kappa) \to
\mathcal{C}(S,\kappa')$ acts by projection $\pi$.
Composition preservation:
$U(d_2 \wedge d_1) = \pi(d_2) \wedge \pi(d_1) =
U(d_2) \wedge U(d_1)$.

\emph{Naturality (first variable).}
For $f: \kappa_1 \to \kappa_2$ witnessed by $d_f$:
\[
\begin{array}{ccc}
  \operatorname{Hom}_\sigma(\kappa_2, \tau)
  & \xrightarrow{\;\Phi_{\kappa_2}\;} & \sigma \\[6pt]
  {\scriptstyle{- \circ F(d_f)}}\downarrow\phantom{x} & &
  \phantom{x}\downarrow{\scriptstyle{\sigma(f)}} \\[6pt]
  \operatorname{Hom}_\sigma(\kappa_1, \tau)
  & \xrightarrow{\;\Phi_{\kappa_1}\;} & \sigma
\end{array}
\]
Top-right: $d_\sigma \xmapsto{\Phi} s \xmapsto{\sigma(f)}
s \wedge \pi_\sigma(d_f)$.
Left-bottom: $d_\sigma \xmapsto{- \circ F(d_f)}
d_\sigma \wedge \iota(\pi_\sigma(d_f)) \xmapsto{\Phi}
s \wedge \pi_\sigma(d_f)$.
Both yield $s \wedge \pi_\sigma(d_f)$. $\checkmark$

\emph{Naturality (second variable).}
For $g: \tau_1 \to \tau_2$ witnessed by $d_g$:
\[
\begin{array}{ccc}
  \operatorname{Hom}_\sigma(\kappa, \tau_1)
  & \xrightarrow{\;\Phi\;} & \sigma \\[6pt]
  {\scriptstyle{U(d_g) \circ -}}\downarrow\phantom{x} & &
  \phantom{x}\downarrow{\scriptstyle{\sigma(g)}} \\[6pt]
  \operatorname{Hom}_\sigma(\kappa, \tau_2)
  & \xrightarrow{\;\Phi\;} & \sigma
\end{array}
\]
Top-right: $d_\sigma \xmapsto{\Phi} s
\xmapsto{\sigma(g)} s \wedge \pi_\sigma(\pi_\tau(d_g))$.
Left-bottom: $d_\sigma \xmapsto{U(d_g) \circ -}
\pi_\tau(d_g) \wedge d_\sigma \xmapsto{\Phi}
s \wedge \pi_\sigma(\pi_\tau(d_g))$.
Both agree. $\checkmark$

\emph{Unit $\eta$:}
$\eta_A: A \to UF(A) = \pi(\iota(A))$.
$\kappa'$ is a direct summand of $\kappa$:
$V = \kappa' \oplus W$. Therefore
$\pi \circ \iota = \mathrm{id}_{\kappa'}$, and
$\eta_A = \mathrm{id}_A$.

\emph{Counit $\varepsilon$:}
$\varepsilon_B: FU(B) = \iota(\pi(B)) \to B$.
On the image of $\iota$, $\varepsilon$ acts as
the identity.

\emph{Zig-zag (i):}
$FUF(A) = \iota(\pi(\iota(A))) = \iota(A) = F(A)$.
$\varepsilon_{F(A)} = \mathrm{id}$.
$(\varepsilon F) \circ (F \eta) = \mathrm{id} \circ
\mathrm{id} = \mathrm{id}_F$. $\checkmark$

\emph{Zig-zag (ii):}
$UFU(B) = \pi(\iota(\pi(B))) = \pi(B) = U(B)$.
$U\varepsilon_B = \mathrm{id}$.
$(U\varepsilon) \circ (\eta U) = \mathrm{id} \circ
\mathrm{id} = \mathrm{id}_U$. $\checkmark$

%----------------------------------------------------------------------
\subsection{Russell dissolution: bilinear form detail}
\label{app:russell-detail}
%----------------------------------------------------------------------

\emph{The profile pairing.}
$B(x, y) = \tau_x(y)$ for $x, y \in V$. Profile
linearity gives linearity in the first argument.
Evaluation linearity (Definition \ref{def:eval-linear})
gives linearity in the second. Therefore $B$ is bilinear.

\emph{Antisymmetrization.}
$\widetilde{B}(x, y) = B(x,y) + B(y,x) = \tau_x(y) +
\tau_y(x)$. On the diagonal:
$\widetilde{B}(x, x) = \tau_x(x) + \tau_x(x) = 0$.
This is the algebraic content of $x \wedge x = 0$.

\emph{Diagonal linearity.}
$Q(x) = B(x, x)$.
$Q(x + y) = B(x+y, x+y) = Q(x) + \widetilde{B}(x,y) +
Q(y)$. Over $\GF(2)$, the Frobenius map $x \mapsto x^2$
is the identity. If $\widetilde{B} = 0$ (symmetric $B$),
then $Q(x+y) = Q(x) + Q(y)$: the diagonal is linear.

\emph{Affine offset.}
$\tau_{d_P}(x) = 1 + Q(x)$.
$\tau_{d_P}(0) = 1 + Q(0) = 1 + 0 = 1 \neq 0$.
Linear maps send $0 \mapsto 0$. Therefore $\tau_{d_P}$
is affine, not linear, and violates profile linearity.
$\checkmark$

%----------------------------------------------------------------------
\subsection{Composition as a 2-cell}
\label{app:composition-2cell}
%----------------------------------------------------------------------

Composition of 1-morphisms $f: A \to B$ (by $d_1$) and
$g: B \to C$ (by $d_2$) produces
$d_1 \wedge d_2 \in \bigwedge^2 V$: a 2-cell.

\emph{Horizontal associativity.}
$(d_1 \wedge d_2) \wedge d_3 = d_1 \wedge (d_2 \wedge d_3)$
by associativity of $\wedge$. $\checkmark$

\emph{Vertical composition.}
2-cells $\alpha: f \Rightarrow f'$ and
$\beta: f' \Rightarrow f''$: $\alpha = d_f \wedge d_{f'}$,
$\beta = d_{f'} \wedge d_{f''}$. Vertical composite
involves $d_{f'} \wedge d_{f'} = 0$, producing
$d_f \wedge d_{f''}$. The intermediate morphism cancels.
$\checkmark$

\emph{1-categorical shadow.}
Identify parallel 1-morphisms connected by an invertible
2-cell. Over $\GF(2)$, all 2-cells are symmetric
(Theorem \ref{thm:groupoid}), hence invertible. The
quotient is a 1-category.

%----------------------------------------------------------------------
\subsection{Power object universal property}
\label{app:power-object}
%----------------------------------------------------------------------

$\mathcal{P}(B) = \Omega^B = [B, \Omega]$ (by Theorem
\ref{thm:closure}). Membership:
${\in_B}: B \times \Omega^B \to \Omega$,
$(b, \phi) \mapsto \phi(b)$.

Given $R \hookrightarrow A \times B$, define
$\chi_R(a)(b) = (\tau_{d_R}(a,b), \sigma_{d_R}(a,b))$.

\emph{Existence:} well-defined by $d_R$'s profile.

\emph{Pullback:}
$(\mathrm{id}_B \times \chi_R)^*(\in_B)(b,a)
= \chi_R(a)(b) = (\tau_{d_R}(a,b), \sigma_{d_R}(a,b))$.
This recovers the full four-valued relation $R$.
$\checkmark$

\emph{Uniqueness:}
Any $\chi'$ pulling back $\in_B$ to $R$ must satisfy
$\chi'(a)(b) = (\tau_{d_R}(a,b), \sigma_{d_R}(a,b))$
for all $(a,b)$. Therefore $\chi' = \chi_R$. $\checkmark$

\emph{Scope limitation:}
The universal property is verified for subobjects defined
by a single discriminator. A full power object must
classify all subobjects, including those defined by
Boolean combinations of discriminators. Extending the
verification to composite subobjects requires showing
that $\chi_R$ is well-defined when $R$ is a
$\kappa$-deformation (not just a single-discriminator
restriction). This follows if $\Omega^B$ contains all
$\GF(2)^2$-valued functions on $B$ (not just those
arising from single discriminators), which holds when
$\kappa$ is sufficiently evolved. The gap: for a fixed
finite $\kappa$, the power object may not classify all
subobjects --- only those witnessable within $\kappa$.
The fibered topos construction resolves this by
computing power objects at the join of the relevant
stages.

\emph{Representability:}
$\Omega^B$ is an object of $\mathcal{C}(S,\kappa)$
because the internal hom $[B, \Omega]$ exists by monoidal
closure. $\checkmark$

%----------------------------------------------------------------------
\subsection{Fano gluing mechanics}
\label{app:fano-gluing}
%----------------------------------------------------------------------

The seven lines of $PG(2, \GF(2))$:
\begin{center}
\small
\begin{tabular}{l|l}
Line & Identity \\
\hline
$\{(1,0,0), (0,1,0), (1,1,0)\}$ &
  $\kappa = \sigma + \tau$ \\
$\{(1,0,0), (0,0,1), (1,0,1)\}$ &
  $\tau$-coverage \\
$\{(0,1,0), (0,0,1), (0,1,1)\}$ &
  $\sigma$-coverage \\
$\{(1,1,0), (0,0,1), (1,1,1)\}$ &
  $\kappa$-coverage \\
$\{(1,0,1), (0,1,0), (1,1,1)\}$ &
  mixed \\
$\{(0,1,1), (1,0,0), (1,1,1)\}$ &
  mixed \\
$\{(1,0,1), (0,1,1), (1,1,0)\}$ &
  mixed
\end{tabular}
\end{center}

Inference: if two points on a line have $\gamma = 1$,
the third is determined by $p_3 = p_1 + p_2$.
After three linearly independent interrogations, all
seven points are covered. $\checkmark$

%----------------------------------------------------------------------
\subsection{Banach--Tarski: group-theoretic analysis}
\label{app:banach-tarski}
%----------------------------------------------------------------------

$SO(3)$ contains a free subgroup $F_2$ on two generators
$a, b$ (rotations with algebraically independent axes).
$F_2$ acts on $S^2$ with orbits that partition the sphere.

Discrimination analysis:
\begin{itemize}[nosep]
  \item \textbf{Group dimension $d_G$}: classifies by
    $F_2$-orbit. Boolean (each point in exactly one orbit).
  \item \textbf{Measure dimension $d_\mu$}: classifies by
    Lebesgue measure. Boolean for measurable sets but
    \emph{undefined} ($\gamma = 0$) for the non-measurable
    pieces from the $F_2$ decomposition.
\end{itemize}

The decomposition is valid on $d_G$ (Boolean, Choice
applies). The ``two spheres of the same volume'' claim
requires $d_\mu$, where Booleanness has not been earned
on the decomposition pieces. The paradox is a type error:
measure-dimension conclusion from group-dimension
construction. $\checkmark$

%----------------------------------------------------------------------
\subsection{Cayley--Dickson step 2: multiplication table}
\label{app:cd-table}
%----------------------------------------------------------------------

$\GF(2)^4 = (\GF(2)^2)^2$ with CD multiplication
$(a,b)(c,d) = (ac + d^*b, da + bc^*)$,
conjugation $(x,y)^* = (y,x)$.

Basis: $e_1 = ((1,0),(0,0))$, $e_2 = ((0,1),(0,0))$,
$e_3 = ((0,0),(1,0))$, $e_4 = ((0,0),(0,1))$.

Sample computation showing non-commutativity:

$e_1 \cdot e_3$: $a=(1,0), b=(0,0), c=(0,0), d=(1,0)$.
$ac + d^*b = (0,0) + (0,1)(0,0) = (0,0)$.
$da + bc^* = (1,0)(1,0) + (0,0) = (1,0)$.
Result: $((0,0),(1,0)) = e_3$.

$e_3 \cdot e_1$: $a=(0,0), b=(1,0), c=(1,0), d=(0,0)$.
$ac + d^*b = (0,0) + (0,0) = (0,0)$.
$da + bc^* = (0,0) + (1,0)(0,1) = (0,0)$.
Result: $0$.

$e_1 \cdot e_3 = e_3 \neq 0 = e_3 \cdot e_1$.
Non-commutativity confirmed. Zero divisors present.
$\checkmark$

%----------------------------------------------------------------------
\subsection{Self-enrichment: full verification}
\label{app:self-enrichment-full}
%----------------------------------------------------------------------

\emph{Hom-objects.}
$[A,B] = \{d \in \kappa : \tau_d(A) = 1,\,\sigma_d(B) = 1\}$
(Definition \ref{def:internal-hom}). Under the
$\sigma$-perspective, this set is a sub-population of
$\kappa$, hence a legitimate object.

\emph{Composition.}
Given $d_1 \in [A,B]$ and $d_2 \in [B,C]$, sequential
filtration produces $d_2 \circ d_1 \in [A,C]$
(Definition \ref{def:composition}). The map
$(d_2, d_1) \mapsto d_2 \circ d_1$ is a $\tau$-discriminator
separating compatible from incompatible pairs.

\emph{Identity.}
$j_A: I \to [A,A]$ sends the grade-0 unit to
$\mathrm{id}_A = 0 \in [A,A]$.

\emph{Associativity}: $(d_3 \circ d_2) \circ d_1 =
d_3 \circ (d_2 \circ d_1)$ by associativity of $\wedge$.
\emph{Unitality}: composing with the degenerate
discriminator leaves any discriminator unchanged.
$\checkmark$

%----------------------------------------------------------------------
\subsection{Subobject classifier: full verification}
\label{app:subobj-full}
%----------------------------------------------------------------------

The ``true'' morphism selects $(1,0) \in \Omega$.
Given $S' \hookrightarrow A$ defined by discriminator $d$
with $S' = \{a : \tau_d(a) = 1,\, \sigma_d(a) = 0\}$:

\emph{Existence}: $\chi_m(a) = (\tau_d(a), \sigma_d(a))$.

\emph{Uniqueness}: any $\chi$ pulling back $\top$ to $S'$
must send $S'$-elements to $(1,0)$ and non-$S'$-elements
to other values. The values are determined by $d$'s profile.

\emph{Pullback}: $S' = \chi_m^{-1}(\{(1,0)\})$. $\checkmark$

%----------------------------------------------------------------------
\subsection{Proof theory: full specification}
\label{app:proof-theory-full}
%----------------------------------------------------------------------

\emph{Connectives on $\Omega = \GF(2)^2$.}
\begin{itemize}[nosep]
  \item Conjunction $\wedge$: componentwise multiplication.
    $(1,0) \wedge (0,1) = (0,0)$: true $\wedge$ false = gap.
  \item Disjunction $\vee$: componentwise maximum.
    $(1,0) \vee (0,1) = (1,1)$: true $\vee$ false = overlap.
  \item Negation $\neg(\tau,\sigma) = (\sigma,\tau)$. Swap.
    $\neg(1,0) = (0,1)$: $\neg$true = false. $\neg(0,0) = (0,0)$:
    $\neg$gap = gap. $\neg(1,1) = (1,1)$: $\neg$overlap = overlap.
  \item Implication $\Rightarrow$: determined by the closure
    adjunction (Theorem \ref{thm:closure}).
\end{itemize}

\emph{Key failures.}
Excluded middle: $p \vee \neg p = (\tau,\sigma) \vee
(\sigma,\tau) = (\max(\tau,\sigma), \max(\sigma,\tau))$.
For $(0,0)$: $(0,0) \vee (0,0) = (0,0) \neq (1,0)$.
Excluded middle fails at the gap.

Explosion: $p \wedge \neg p = (\tau \cdot \sigma,
\sigma \cdot \tau)$. For $(1,1)$:
$(1,1) \wedge (1,1) = (1,1)$, not $(1,0)$. The overlap
does not entail arbitrary conclusions.

\emph{Relationship to logical traditions.}
Classical logic: the Boolean restriction to
$\{(1,0), (0,1)\}$. Excluded middle holds. Local
achievement.

Intuitionistic logic: shares failure of excluded middle
but for a different reason. Intuitionistically,
$p \vee \neg p$ fails because disjunction requires a
witness. Here it fails because neither predicate fires
(epistemic, not constructive).

Paraconsistent logic: the overlap $(1,1)$ tolerates
contradiction without explosion. The logic is
paraconsistent at the overlap.

\emph{Dynamic layer.}
$\kappa$-evolution resolves residues: elements at
$(0,0)$ or $(1,1)$ are moved toward $\{(1,0), (0,1)\}$
by introducing finer discriminators. Each step
monotonically increases information (Proposition
\ref{prop:monotone}).

%----------------------------------------------------------------------
\subsection{Self-hosting: full proof}
\label{app:self-hosting-full}
%----------------------------------------------------------------------

The Boolean dependent type restricts $\Omega$ to
$\{(1,0), (0,1)\}$. Within this restriction:
\begin{itemize}[nosep]
  \item Two elements: $(1,0)$ and $(0,1)$. Identify
    $(1,0) \leftrightarrow 1$ and $(0,1) \leftrightarrow 0$.
  \item Addition: $(1,0) + (0,1) = (1,1)$, but within the
    Boolean restriction, $(1,1)$ is the overlap --- not a
    Boolean state. The XOR: $1 + 1 = 0$ (mapping back to
    the Boolean restriction). This gives $\GF(2)$ addition.
  \item Multiplication: componentwise multiplication.
    $(1,0) \cdot (1,0) = (1,0)$: $1 \cdot 1 = 1$.
    $(1,0) \cdot (0,1) = (0,0) \to 0$: $1 \cdot 0 = 0$.
\end{itemize}

The internally constructed $\GF(2)$ has the same algebraic
properties as the external one: characteristic 2,
$1 + 1 = 0$, $-x = x$. The construction proceeds:
$\GF(2)_{\mathrm{ext}} \to$ four-valued logic $\to$
Boolean type $\to \GF(2)_{\mathrm{int}} \to \cdots$.
The self-reference is non-pathological:
$\GF(2)_{\mathrm{int}}$ is algebraically identical to
$\GF(2)_{\mathrm{ext}}$, and $x \wedge x = 0$ prevents
contradictory fixpoints. $\checkmark$

%----------------------------------------------------------------------
\subsection{Fixed-point theorem: full proof}
\label{app:fixed-point-full}
%----------------------------------------------------------------------

Three structures on the $\GF(2)$ toroid:

(i) \emph{The adjunction}: the free-forgetful adjunction
(Theorem \ref{thm:adjunction}). A predicate is this
adjunction --- a geometric projection. ``Observation'' is
functor traversal; ``firing'' is profile retention under
projection.

(ii) \emph{$\GF(2)$}: the target of the projections. Over
$\GF(2)$, projection and inclusion are related by
$\pi \circ \iota = \mathrm{id}$ (direct summand), making
the adjunction a retract.

(iii) \emph{Belnap logic}: the $\GF(2) \times \GF(2)$
classifier. This is the product of two copies of the
adjunction's target, giving four-valued truth.

The three are the non-zero points of $\GF(2)^2$. The
2-cell filling their triangle is the operation mediating
between them: the wedge product that composes projections
into a coherent discrimination. This is discrimination.

The three coexist (each pair determines the third via the
triangle identity), none requires an external observer,
and the 2-cell is the unique stable completion of the
toroid's geometry. $\checkmark$

%----------------------------------------------------------------------
\subsection{Irremovability: full proof}
\label{app:irremovable-full}
%----------------------------------------------------------------------

Assume the 2-cell (discrimination) is absent. Four
consequences:

\emph{1. Sheaf failure.} Without the 2-cell, covering
morphisms (the six adjunctions) don't exist. Local data
on $\kappa, \sigma, \tau$ can't be assembled into global
sections. The site $\mathfrak{S}$ loses its Grothendieck
topology. The topos degrades to three disconnected
presheaves.

\emph{2. Manifold failure.} The toroid's atlas consists of
charts related by the six adjunctions as transition maps.
Without the 2-cell, the transition maps vanish. Three
isolated points, no manifold.

\emph{3. Bimodule failure.} $\sigma$ as a bimodule between
$\kappa$ and $\tau$ loses its mediating structure.
$\kappa$ and $\tau$ become algebraically isolated.

\emph{4. Self-resolution.} The unfilled 2-cell is a
homological gap (non-trivial $H_1$ class). By
exhaustivity (Theorem \ref{thm:exhaustive}), the gap
specifies its own filling. The specified filling has
boundary $\{$adjunction, $\GF(2)$, Belnap$\}$. That
filling is discrimination. Removal triggers reinstatement.

This is internal necessity: within the framework's
machinery, discrimination cannot be absent because
the absence is a gap whose filling is discrimination.
A reader who has not accepted the operationalist axiom
may reject the framework's gap-filling machinery, in
which case the argument has no external force.
$\checkmark$

%----------------------------------------------------------------------
\subsection{Sheaf topos: full proof}
\label{app:sheaf-full}
%----------------------------------------------------------------------

The site $\mathfrak{S}$ has $\kappa$-states as objects and
sub-regime inclusions as morphisms. The Grothendieck
topology: a family $\{\kappa_i \hookrightarrow \kappa\}$
covers $\kappa$ if the $\kappa_i$ jointly contain all
discriminators of $\kappa$.

The six triadic adjunctions provide covering data: the
three axes ($\kappa$, $\sigma$, $\tau$) with two chiralities
each give six perspectives, any three linearly independent
ones covering the full regime (Theorem \ref{thm:convergence}).

Sheaf condition: a presheaf $F$ on $\mathfrak{S}$ is a
sheaf iff for every covering family, compatible local
sections glue to a unique global section. Compatibility
is order-independence (Proposition \ref{prop:monoidal}):
the $\tau/\sigma$ profile doesn't depend on which covering
morphism was applied. Gluing is the assembly of local
profiles into a global profile.

The presheaf-to-sheaf transition is $\kappa$-evolution:
a presheaf that fails the gluing condition has a residue
(the incompatibility), which triggers evolution adding the
discriminator needed to resolve it. $\checkmark$

%----------------------------------------------------------------------
\subsection{FOL derivation: full construction}
\label{app:fol}
%----------------------------------------------------------------------

\emph{Universal quantifier.}
$\forall_f: \Omega^A \to \Omega^B$ (right adjoint to
pullback $f^*$): $(\forall_f P)(b) = \bigwedge_{a \in
f^{-1}(b)} P(a)$, where $\bigwedge$ is componentwise
multiplication in $\GF(2)^2$. The adjunction
$f^* \dashv \forall_f$ holds because $f^*(Q) \leq P$ iff
$Q \leq \forall_f P$.

\emph{Existential quantifier.}
$\exists_f: \Omega^A \to \Omega^B$ (left adjoint to $f^*$):
$(\exists_f P)(b) = \bigvee_{a \in f^{-1}(b)} P(a)$,
componentwise maximum.

\emph{Modus ponens.} Evaluation morphism
$\mathrm{ev}: [\Omega, \Omega] \times \Omega \to \Omega$:
$(h, p) \mapsto h(p)$. A morphism in the topos.

\emph{Universal instantiation.} Counit of
$\forall_f \dashv f^*$: maps $(\forall_f P)(f(a))$ to
$P(a)$.

\emph{Existential generalization.} Unit of
$f^* \dashv \exists_f$: maps $P(a)$ to
$(\exists_f P)(f(a)) \geq P(a)$.

Terms = morphisms. Variables = elements. Substitution =
composition. All internal. $\checkmark$

%----------------------------------------------------------------------
% Missing labels referenced from other sections
%----------------------------------------------------------------------

% prop:presheaf-sheaf is referenced from sec04
% We provide it as a cross-reference target
\begin{proposition}[Presheaf-to-sheaf transition]
\label{prop:presheaf-sheaf}
A presheaf on the site $\mathfrak{S}$ that fails the
gluing condition has a residue (the incompatible local
sections). This residue triggers $\kappa$-evolution,
adding the discriminator needed to resolve the
incompatibility. The presheaf becomes a sheaf after
finitely many evolutions (bounded by $3$ per dimension,
Theorem \ref{thm:convergence}).
\end{proposition}

% prop:commutative is referenced from sec08
\begin{proposition}[Commutativity of $\wedge$ over $\GF(2)$]
\label{prop:commutative}
The wedge product over $\GF(2)$ is commutative:
$a \wedge b = b \wedge a$. This follows from
$a \wedge b = -b \wedge a$ (antisymmetry) and $-1 = 1$
in $\GF(2)$.
\end{proposition}

%----------------------------------------------------------------------
\subsection{Foundation: the partial result}
\label{app:foundation-full}
%----------------------------------------------------------------------

Self-membership ($a \in a$) is excluded by profile
linearity (Theorem \ref{thm:russell}): the self-membership
profile is affine, not linear.

Cycles of length $\geq 2$ are algebraically permitted.
A cycle $a \in b \in a$ is witnessed by independent
$d_1, d_2$ with $\tau_{d_1}(a) = 1, \sigma_{d_1}(b) = 1,
\tau_{d_2}(b) = 1, \sigma_{d_2}(a) = 1$. No algebraic
obstruction: $d_1 \wedge d_2 \neq 0$, no linearity
violation.

Infinite chains are prevented by finiteness of $\kappa$:
with $\dim(\kappa) = n$, at most $2^n - 1$ non-zero
discriminators exist, bounding the number of witnessable
links. $\checkmark$

%----------------------------------------------------------------------
\subsection{Choice: finite and infinite cases}
\label{app:choice-full}
%----------------------------------------------------------------------

\emph{Finite case.} Constructive: each $S_i$ contains
finitely many $\kappa$-equivalence classes. Selection
failure produces a residue (Theorem \ref{thm:exhaustive})
specifying the needed discriminator. Iteration terminates
in $\leq 3$ evolutions per dimension. $\checkmark$

\emph{Infinite case.} Application of the operationalist
axiom (Definition \ref{def:operationalist}): no
discriminator separates ``completed selection function''
from ``constructively obtainable for every finite
sub-family.'' This is an axiom application, not an
algebraic derivation. A Platonist may reject the
axiom. $\checkmark$

%----------------------------------------------------------------------
\subsection{Exhaustivity: the $\partial \circ \partial = 0$
  argument}
\label{app:exhaustivity-full}
%----------------------------------------------------------------------

A residue $R_d(S)$ occupies a cycle in $\ker \partial_k$
not in $\operatorname{im} \partial_{k+1}$: a non-trivial
homology class $[R] \in H_k$.

$\partial \circ \partial = 0$ means the boundary of
a boundary is empty: the residue of a residue is zero.
The gap is fully characterized at detection.

The class $[R]$ specifies: the grade $k$, the boundary
$(k{-}1)$-cells, and (when $H_k(\kappa) = 0$) a unique
filling.

\emph{Constructive filling procedure.}
Let $R \in \ker \partial_k$ be a $k$-cycle (a residue).
Introduce a new basis discriminator $e_{n+1}$ independent
of $\kappa$ (a $\kappa$-evolution). Form the $(k{+}1)$-chain
$C = R \wedge e_{n+1}$.

By the Leibniz rule over $\GF(2)$
($\partial(a \wedge b) = \partial a \wedge b + a \wedge
\partial b$, with $(-1)^{|a|} = 1$ since $-1 = 1$):
\[
  \partial(R \wedge e_{n+1}) = \partial R \wedge e_{n+1}
  + R \wedge \partial(e_{n+1}) = 0 \cdot e_{n+1}
  + R \cdot 1 = R.
\]
Here $\partial R = 0$ because $R$ is a cycle, and
$\partial(e_{n+1}) = 1$ (the grade-1 boundary is the
augmentation: removing the single factor yields the
grade-0 unit).

Therefore $\partial C = R$: the residue becomes a boundary
in the extended complex $\bigwedge(\kappa \cup \{e_{n+1}\})$.
After $\kappa$-evolution, $[R] = 0$ in $H_k(\kappa')$.

The procedure is: \emph{$\kappa$-evolution itself is the
filling}. Any independent discriminator works. The residue
$R$ specifies \emph{that} a filling exists (by its shape
as a cycle) and \emph{what} the filling looks like (wedge
with any independent direction). From the inner system's
perspective, the filling doesn't exist (it requires a
discriminator not in $\kappa$). From the outer system's
perspective, it is trivially constructible. The gap
between inner and outer is exactly one discriminator.
$\checkmark$

%----------------------------------------------------------------------
\subsection{Groupoid structure}
\label{app:groupoid-full}
%----------------------------------------------------------------------

Over $\GF(2)$, $d \wedge d' = d' \wedge d$ (commutativity,
Proposition \ref{prop:commutative}). Every 2-cell
$d_1 \wedge d_2$ equals $d_2 \wedge d_1$: the homotopy
from $d_1$ to $d_2$ is the same as the homotopy from
$d_2$ to $d_1$. Every 2-morphism is its own inverse.
The 2-morphisms form a symmetric groupoid. By induction,
$k$-morphisms for $k \geq 2$ are all invertible:
the homotopy tower is a groupoid at every level above
grade 1. $\checkmark$

%----------------------------------------------------------------------
\subsection{2-category axioms: full check}
\label{app:2cat-full}
%----------------------------------------------------------------------

\emph{Horizontal associativity}: $(d_1 \wedge d_2) \wedge
d_3 = d_1 \wedge (d_2 \wedge d_3)$ by $\wedge$.

\emph{Vertical composition}: 2-cells
$\alpha = d_f \wedge d_{f'}$, $\beta = d_{f'} \wedge d_{f''}$.
Vertical composite: $(d_{f'} \wedge d_{f''}) \wedge
(d_f \wedge d_{f'})$ contains $d_{f'} \wedge d_{f'} = 0$,
producing $d_f \wedge d_{f''}$. Intermediate cancels.

\emph{Identity 2-cell}: $\mathrm{id}_f = d_f \wedge d_f = 0$
(the degenerate 2-cell).

\emph{Interchange law.}
The verification works in the $\GF(2)^2$ profile space,
not in $\bigwedge^4 V$. Each morphism $d$ has profile
$p(d)(a) = (\tau_d(a), \sigma_d(a)) \in \GF(2)^2$.
Each 2-cell $\alpha: f \Rightarrow f'$ has profile
$p(\alpha) = p(f) + p(f')$. Composition is componentwise
multiplication in $\GF(2)^2$.

\emph{Path 1}: $(\beta * f') \circ (g * \alpha)$.
$g * \alpha = p(g) \cdot (p(f) + p(f')) = g \cdot f +
g \cdot f'$. $\beta * f' = (p(g) + p(g')) \cdot p(f') =
g \cdot f' + g' \cdot f'$. Vertical composite:
$(g \cdot f + g \cdot f') + (g \cdot f' + g' \cdot f')
= g \cdot f + 2 \cdot g \cdot f' + g' \cdot f'$. Over
$\GF(2)$: $2 = 0$, so $= g \cdot f + g' \cdot f'$.

\emph{Path 2}: $(g' * \alpha) \circ (\beta * f)$.
$\beta * f = (p(g) + p(g')) \cdot p(f) =
g \cdot f + g' \cdot f$. $g' * \alpha =
p(g') \cdot (p(f) + p(f')) = g' \cdot f + g' \cdot f'$.
Vertical composite:
$(g \cdot f + g' \cdot f) + (g' \cdot f + g' \cdot f')
= g \cdot f + 2 \cdot g' \cdot f + g' \cdot f'$. Over
$\GF(2)$: $= g \cdot f + g' \cdot f'$.

Both paths yield $p(g \circ f) + p(g' \circ f')
\in \GF(2)^2$. The cross terms cancel because
$1 + 1 = 0$ in $\GF(2)$. The result is grade 2 in
the profile space (a $\GF(2)^2$ element), not grade 4:
the Klein four-group's characteristic 2 collapses the
interchange square from $\bigwedge^4$ to $\GF(2)^2$.

This is the same $1 + 1 = 0$ that excludes
self-reference ($v \wedge v = 0$), kills the Russell
predicate's affine offset, and makes the halting
tower level-invariant. The interchange law holds for
the same reason everything else in the framework
holds: characteristic 2. $\checkmark$

%----------------------------------------------------------------------
\subsection{Yoneda}
\label{app:yoneda-full}
%----------------------------------------------------------------------

The representable presheaf $h_A: \mathcal{C}^{op} \to
\mathbf{Set}$ sends $B \mapsto \operatorname{Hom}(B, A) =
\{d \in \kappa : \tau_d(B) = 1, \sigma_d(A) = 1\}$.

Yoneda: $\operatorname{Nat}(h_A, F) \cong F(A)$ for any
presheaf $F$. Proof: a natural transformation
$\alpha: h_A \Rightarrow F$ is determined by
$\alpha_A(\mathrm{id}_A) \in F(A)$ (the image of the
identity). Naturality forces all other components:
$\alpha_B(f) = F(f)(\alpha_A(\mathrm{id}_A))$. $\checkmark$

%----------------------------------------------------------------------
\subsection{Delooping and monoidal closure}
\label{app:deloop-closure}
%----------------------------------------------------------------------

\emph{Symmetric monoidal strictness.}
The monoidal product $\otimes = \wedge$ is \emph{strictly}
associative: $(a \wedge b) \wedge c = a \wedge (b \wedge c)$
as elements of $\bigwedge V$, not merely up to isomorphism.
The associator $\alpha_{A,B,C}$ is the identity. The unit
$I$ is the grade-0 element, with $d \wedge 1 = d = 1 \wedge d$
strictly. The left and right unitors $\lambda, \rho$ are
identities. The symmetry $\beta_{A,B}: a \wedge b = b \wedge a$
is the identity over $\GF(2)$ (since $-1 = 1$).

For a strict symmetric monoidal category, the pentagon,
triangle, and hexagon diagrams commute trivially: every
path is the identity, so all paths agree. No coherence
diagrams need explicit verification.

At the fiber level: the join $\kappa_1 \vee \kappa_2$ is
strictly associative, commutative, and unital ($\kappa \vee
\varnothing = \kappa$). The Galois group reduces under
joins: $\mathrm{Gal}(V/(\kappa_1 \vee \kappa_2)) \subseteq
\mathrm{Gal}(V/\kappa_1) \cap \mathrm{Gal}(V/\kappa_2)$.
Intersection of subgroups is associative and commutative,
mirroring the strict monoidal structure at the object level.
$\checkmark$

\emph{Delooping}: The $\sigma$-perspective views
discriminators as a population. A monoidal $(n{-}1)$-category
with one object (the entire $\sigma$-population) whose
morphisms are $\kappa$-discriminators and monoidal product
is $\wedge$. This is the standard delooping: a 2-category
with one object is a monoidal 1-category.

\emph{Monoidal closure}: $\operatorname{Hom}(A \otimes B, C)
\cong \operatorname{Hom}(A, [B, C])$ where
$[B, C] = \{d \in \kappa : \tau_d(B) = 1, \sigma_d(C) = 1\}$.
The bijection: a morphism $A \otimes B \to C$ witnessed by
$d$ with $\tau_d(A \otimes B) = 1, \sigma_d(C) = 1$
decomposes (by the monoidal product structure) into a
morphism $A \to [B, C]$ witnessed by the same $d$ restricted
to the $A$-component. Naturality in $A$: for $h: A' \to A$ witnessed by $d_h$,
precomposition with $h \otimes \mathrm{id}_B = d_h \wedge 1$
on the left of the bijection corresponds to precomposition
with $d_h$ on the right. Since $d_h \wedge 1 = d_h$
(strict unit), the naturality square commutes:
restricting $d_h \wedge 1$ to the $A$-component yields
$d_h$. At the fiber level, this is functoriality of the
restriction map along $\kappa$-inclusions, equivariant
under Galois group reduction. $\checkmark$

Naturality also follows from $\pi$ being
a ring homomorphism. $\checkmark$

%----------------------------------------------------------------------
\subsection{Boolean dependent type and convergence rate}
\label{app:bool-convergence}
%----------------------------------------------------------------------

\emph{Boolean type}: $\chi_d(a) = \tau_d(a) + \sigma_d(a)
= 1$ for all $a \in S$ is a predicate on $d$ (not on
elements). It depends on the specific discriminator and
population. Classical logic holds iff this predicate is
satisfied. The type is dependent: its validity is a
function of the dimension, not a global property.

\emph{Convergence rate}: the Fano plane has 7 points and 7
lines, each line through 3 points. Three non-collinear
points determine all 7 via Fano inference ($p_3 = p_1 +
p_2$ on each line). Three linearly independent
interrogations are non-collinear (if they were collinear,
one would be the sum of the other two, contradicting
independence). Therefore 3 suffices. $\checkmark$

%----------------------------------------------------------------------
\subsection{Discriminator semantics: single predicate,
  composed pairs, dependent types}
\label{app:discriminator-semantics}
%----------------------------------------------------------------------

The foundational definitions (\S\ref{sec:framework})
rest on three distinctions that must not be conflated:

\emph{1. A discriminator is a single predicate.}
A discriminator $d \in \bigwedge^1 V$ provides one map
$d: S \to \GF(2)$: fires or silent. It does not produce
a pair. It tests for one property.

\emph{2. $\tau$ and $\sigma$ are datum properties, not
discriminator outputs.}
A datum $a$ may carry the $\tau$ label, the $\sigma$
label, both, or neither. These are attributes of the
data. A $\tau$-keyed discriminator $d_\tau$ checks
whether $\tau$ is active: $d_\tau(a) = 1$ means $\tau$
is active for $a$. A $\sigma$-keyed discriminator
$d_\sigma$ checks $\sigma$ independently.

\emph{3. The membership profile is a composed
discriminator.}
The pair $(\tau(a), \sigma(a))$ is the result of
composing two base discriminators: $d_\tau$ and
$d_\sigma$. This composition corresponds to the
grade-2 element $d_\tau \wedge d_\sigma \in
\bigwedge^2 V$. The wedge product is not a formal
algebraic operation applied after the fact --- it
\emph{is} the comparison of two discriminator results.
The grade-2 element records the joint outcome.

If $d_\tau = d_\sigma$ (the same discriminator used
for both properties), then
$d_\tau \wedge d_\sigma = d \wedge d = 0$: self-comparison
is degenerate. A non-trivial profile requires two
\emph{independent} discriminators.

\emph{4. Higher-level discriminators are dependent types.}
A discriminator that tests a composite property ---
such as the XOR signature $\chi(a) = \tau(a) + \sigma(a)$,
which tests whether exactly one base property is
active --- depends on the results of the base
discriminators. Its type is dependent: it presupposes
that $d_\tau$ and $d_\sigma$ have been applied.

The Boolean dimension $\chi(a) = 1$ for all $a$ is the
principal dependent type: it asserts that every datum
carries exactly one of $\tau$ or $\sigma$, which requires
both base discriminators to have been evaluated on every
element.

\emph{5. $S_3$ rotations act on compositions.}
The triangle identity $\kappa = \sigma + \tau$ relates
three elements of $\GF(2)^2$. Under $S_3$ permutation,
the roles change: which discriminator is keyed to which
property, and therefore which composition is being formed.
The six adjunctions of Theorem \ref{thm:adjunction} are
six ways of pairing base discriminators --- six
compositions $d_A \wedge d_B$ where $A, B$ range over
$\{\kappa, \sigma, \tau\}$. Rotating the triangle changes
the composition, not the data. $\checkmark$


%----------------------------------------------------------------------
\subsection{Symmetric monoidal coherence: exhaustive
  enumeration}
\label{app:monoidal-coherence}
%----------------------------------------------------------------------

The monoidal product $\otimes = \wedge$ on
$\mathcal{C}(S,\kappa)$ is strict. We enumerate every
coherence condition and verify each is the identity.

\emph{1. Associator.}
$\alpha_{A,B,C}: (A \otimes B) \otimes C \to A \otimes
(B \otimes C)$.
Computation: $(d_A \wedge d_B) \wedge d_C = d_A \wedge
(d_B \wedge d_C)$ by associativity of $\wedge$ in
$\bigwedge V$. Both sides are the same element of
$\bigwedge^3 V$. $\alpha_{A,B,C} = \mathrm{id}$.
$\checkmark$

\emph{2. Left unitor.}
$\lambda_A: I \otimes A \to A$ where $I = 1 \in
\bigwedge^0 V \cong \GF(2)$.
Computation: $1 \wedge d_A = d_A$ (the grade-0 unit acts
as the identity for $\wedge$). $\lambda_A = \mathrm{id}$.
$\checkmark$

\emph{3. Right unitor.}
$\rho_A: A \otimes I \to A$.
Computation: $d_A \wedge 1 = d_A$. $\rho_A = \mathrm{id}$.
$\checkmark$

\emph{4. Braiding.}
$\beta_{A,B}: A \otimes B \to B \otimes A$.
Computation: $d_A \wedge d_B = (-1)^{|d_A||d_B|}
d_B \wedge d_A = d_B \wedge d_A$ (since $-1 = 1$ in
$\GF(2)$). For grade-1 elements:
$d_A \wedge d_B = d_B \wedge d_A$ literally.
$\beta_{A,B} = \mathrm{id}$. $\checkmark$

\emph{5. Inverse braiding.}
$\beta_{A,B}^{-1} = \beta_{B,A}$. Since
$\beta_{A,B} = \mathrm{id}$, also
$\beta_{A,B}^{-1} = \mathrm{id}$. $\checkmark$

\emph{6. Pentagon diagram.}
The five edges of the pentagon are:
\begin{align*}
  &((A \otimes B) \otimes C) \otimes D
    \xrightarrow{\alpha_{A \otimes B, C, D}}
  (A \otimes B) \otimes (C \otimes D)
    \xrightarrow{\alpha_{A, B, C \otimes D}}
  A \otimes (B \otimes (C \otimes D)) \\
  &((A \otimes B) \otimes C) \otimes D
    \xrightarrow{\alpha_{A,B,C} \otimes \mathrm{id}_D}
  (A \otimes (B \otimes C)) \otimes D
    \xrightarrow{\alpha_{A, B \otimes C, D}}
  A \otimes ((B \otimes C) \otimes D) \\
  &\qquad\qquad\qquad\qquad\qquad
    \xrightarrow{\mathrm{id}_A \otimes \alpha_{B,C,D}}
  A \otimes (B \otimes (C \otimes D))
\end{align*}
Every $\alpha = \mathrm{id}$. Top path:
$\mathrm{id} \circ \mathrm{id} = \mathrm{id}$.
Bottom path:
$\mathrm{id} \circ \mathrm{id} \circ \mathrm{id}
= \mathrm{id}$. Both paths = id. $\checkmark$

\emph{7. Triangle diagram.}
$(A \otimes I) \otimes B
\xrightarrow{\alpha_{A,I,B}}
A \otimes (I \otimes B)
\xrightarrow{\mathrm{id}_A \otimes \lambda_B}
A \otimes B$
versus
$(A \otimes I) \otimes B
\xrightarrow{\rho_A \otimes \mathrm{id}_B}
A \otimes B$.

Top path: $(\mathrm{id}_A \otimes \lambda_B) \circ
\alpha_{A,I,B} = \mathrm{id} \circ \mathrm{id}
= \mathrm{id}$.
Bottom path: $\rho_A \otimes \mathrm{id}_B = \mathrm{id}
\otimes \mathrm{id} = \mathrm{id}$.
Both = id. $\checkmark$

\emph{8. First hexagon diagram.}
$A \otimes (B \otimes C)
\xrightarrow{\beta_{A, B \otimes C}}
(B \otimes C) \otimes A
\xrightarrow{\alpha_{B,C,A}}
B \otimes (C \otimes A)$
versus
$A \otimes (B \otimes C)
\xrightarrow{\alpha_{A,B,C}^{-1}}
(A \otimes B) \otimes C
\xrightarrow{\beta_{A,B} \otimes \mathrm{id}_C}
(B \otimes A) \otimes C
\xrightarrow{\alpha_{B,A,C}}
B \otimes (A \otimes C)
\xrightarrow{\mathrm{id}_B \otimes \beta_{A,C}}
B \otimes (C \otimes A)$.

Every morphism is id. Top: $\mathrm{id}^3 = \mathrm{id}$.
Bottom: $\mathrm{id}^5 = \mathrm{id}$. $\checkmark$

\emph{9. Second hexagon diagram.}
Identical analysis with $\beta^{-1}$ replacing $\beta$.
Since $\beta = \beta^{-1} = \mathrm{id}$, the diagram
is identical to the first. $\checkmark$

\emph{10. Symmetry condition.}
$\beta_{B,A} \circ \beta_{A,B} = \mathrm{id}_{A \otimes B}$.
$\mathrm{id} \circ \mathrm{id} = \mathrm{id}$. $\checkmark$

\medskip
All ten coherence conditions are verified by the same
mechanism: the monoidal structure on
$\mathcal{C}(S,\kappa)$ is \emph{strict symmetric
monoidal} because $\wedge$ is strictly associative,
commutative (over $\GF(2)$), and unital. Every coherence
morphism is the identity. This is not an inductive
argument (``strict implies trivial'') but an exhaustive
constructive enumeration: each of the ten conditions is
individually checked and found to be the identity.


%----------------------------------------------------------------------
\subsection{Finite limits: exhaustive construction}
\label{app:limits-exhaustive}
%----------------------------------------------------------------------

We enumerate each finite limit type in
$\mathcal{C}(S,\kappa)$, constructing the limit object
explicitly and verifying the universal property.
Profile linearity (Definition \ref{def:profile-linear})
makes all constraints linear over $\GF(2)$; universality
follows from uniqueness of linear factorizations.

\emph{1. Terminal object.}
\emph{Construction}: $1 = S$ (the entire population),
classified by the zero discriminator $d = 0$
($\tau_0(a) = 0$, $\sigma_0(a) = 0$ for all $a$: no
element is distinguished). \emph{Cone}: for any object
$X$, the unique morphism $X \to 1$ is witnessed by the
zero discriminator (every element maps to the
undifferentiated whole). \emph{Universality}: the zero
discriminator is the unique discriminator that makes no
distinction, so the morphism is unique. $\checkmark$

\emph{2. Product $A \times B$.}
\emph{Construction}: let $d_A$ classify $A$ and $d_B$
classify $B$. The product is the equivalence class under
the joint profile $(d_A, d_B)$ in the regime
$\kappa' = \kappa \vee \{d_A, d_B\}$. Projections:
$\pi_A$ extracts the $d_A$-component, $\pi_B$ extracts
$d_B$.
\emph{Cone}: given $f: X \to A$ (by $d_f$) and
$g: X \to B$ (by $d_g$), the pairing
$\langle f, g \rangle: X \to A \times B$ is
witnessed by $d_f + d_g$ (sum over $\GF(2)$).
\emph{Universality}: $\pi_A \circ \langle f, g \rangle$
extracts the $d_A$-component of $d_f + d_g$. By profile
linearity, this is $d_f$'s contribution to the
$A$-classification, which is $f$. Similarly for $\pi_B$.
\emph{Uniqueness}: any $h: X \to A \times B$ with
$\pi_A \circ h = f$ and $\pi_B \circ h = g$ must have
profile agreeing with $d_f$ on $A$ and $d_g$ on $B$.
Over $\GF(2)$, a linear map is determined by its values
on a basis; the two projection conditions determine $h$
uniquely. $\checkmark$

\emph{3. Equalizer.}
\emph{Construction}: given $f, g: A \to B$ witnessed by
$d_f, d_g$, the equalizer is $E = \{a \in A :
\tau_{d_f}(a) = \tau_{d_g}(a),\; \sigma_{d_f}(a) =
\sigma_{d_g}(a)\}$ --- the elements where $f$ and $g$
agree. The equalizer discriminator is $d_f + d_g$: it
fires ($= 1$) exactly where $f$ and $g$ differ (since
$1 + 1 = 0$ and $0 + 0 = 0$ in $\GF(2)$, the sum is
zero where they agree and nonzero where they differ).
$E = \ker(d_f + d_g)$.
\emph{Inclusion}: $\iota: E \hookrightarrow A$ is the
sub-regime inclusion.
\emph{Universality}: for any $h: X \to A$ with
$f \circ h = g \circ h$, the image of $h$ lies in
$\ker(d_f + d_g) = E$. The unique factorization
$\bar{h}: X \to E$ is $h$ with codomain restricted to
$E$. Uniqueness: $\iota$ is injective (a monomorphism),
so the factorization is unique. $\checkmark$

\emph{4. Pullback.}
\emph{Construction}: given $f: A \to C$ (by $d_f$) and
$g: B \to C$ (by $d_g$), the pullback is
$P = \{(a, b) \in A \times B : f(a) = g(b)\}$.
In the exterior algebra: $P$ is the kernel of
$d_f + d_g$ restricted to $A \times B$ (elements where
the $C$-profiles agree). The pullback discriminator is
$d_f + d_g$ evaluated on the joint $(A, B)$ population.
\emph{Projections}: $p_A: P \to A$ and $p_B: P \to B$
extract the components.
\emph{Universality}: given $h: X \to A$ and $k: X \to B$
with $f \circ h = g \circ k$, the image of
$(h, k): X \to A \times B$ lies in $\ker(d_f + d_g) = P$.
The unique factorization is $(h, k)$ restricted to $P$.
\emph{Uniqueness}: determined by the two projection
conditions, which fix $(h, k)$ uniquely (as in the
product case). $\checkmark$

\emph{5. General finite limit.}
A finite limit is a finite combination of products and
equalizers (or equivalently, a terminal object and
pullbacks). Since each component is constructed above,
any finite limit is constructed by iterating these
operations. At each step, the construction adds linear
constraints over $\GF(2)$. The solution set is a
linear subspace (intersection of kernels). The
universal property follows from linearity: the
factoring morphism is the unique linear map satisfying
the projection conditions.

If the constraints require discriminators not in $\kappa$,
$\kappa$-evolution provides them (Theorem
\ref{thm:exhaustive}, the constructive filling procedure
of \S\ref{app:exhaustivity-full}): the residue of
incompatible constraints is a cycle $R$, and
$R \wedge e_{n+1}$ fills it. After evolution, the limit
exists in $\kappa'$.

The fibered topos computes finite limits at the join
$\kappa_\vee$ of all stages involved (Definition
\ref{def:fibered}), ensuring all needed discriminators
are present. $\checkmark$


%----------------------------------------------------------------------
\subsection{Power objects: composite subobjects at joins}
\label{app:power-composite}
%----------------------------------------------------------------------

We extend the single-discriminator verification (B.4) to
composite subobjects, using the Galois structure of the
$\kappa$-split.

\emph{1. Single-discriminator subobjects} (verified in B.4).
A subobject of $B$ defined by $d \in \kappa$ is classified
by $\chi_d: A \to \Omega$ with $\chi_d(a)(b) =
(\tau_d(a,b), \sigma_d(a,b))$. Existence, pullback, and
uniqueness checked. $\checkmark$

\emph{2. Linear-combination subobjects.}
A subobject defined by $d_1 + d_2$ (sum in $V$) is a
valid discriminator by profile linearity
(Definition \ref{def:profile-linear}):
$(d_1 + d_2)(a) = d_1(a) + d_2(a)$.
Its classifying morphism is
$\chi_{d_1 + d_2}(a)(b) = (\tau_{d_1+d_2}(a,b),
\sigma_{d_1+d_2}(a,b))$, which by profile linearity
equals the componentwise sum of the individual
classifications. This is a well-defined element of
$\Omega^B$, unique by the same argument as case 1.
$\checkmark$

\emph{3. Wedge-product subobjects.}
A subobject defined by a grade-2 discriminator
$d_1 \wedge d_2 \in \bigwedge^2 V$ classifies elements
by the \emph{joint} profile of $d_1$ and $d_2$.
At stage $\kappa$ containing both $d_1$ and $d_2$:
$\chi_{d_1 \wedge d_2}(a)(b) = (d_1(a,b) \cdot d_2(a,b),
\; \ldots)$ where the product is componentwise in
$\GF(2)^2$.

The classifying morphism is well-defined because
$d_1, d_2 \in \kappa$ and evaluation is
$\GF(2)$-valued. Uniqueness: the pullback condition
determines each component. $\checkmark$

\emph{4. General subobjects at stage $\kappa$.}
Any subobject of $B$ classifiable at stage $\kappa$ is
defined by a collection of discriminators $\{d_i\}
\subset \kappa$ and their algebraic combinations
(sums and wedge products). By profile linearity, the
classifiable subobjects form a \emph{linear subspace}
of $\operatorname{Hom}(B, \Omega)$. The power object
$\Omega^B$ at stage $\kappa$ is this linear subspace.
The universal property holds because:
\begin{itemize}[nosep]
  \item \emph{Existence}: for any relation
    $R \hookrightarrow A \times B$ defined by
    discriminators in $\kappa$, $\chi_R$ is constructed
    from those discriminators.
  \item \emph{Uniqueness}: $\chi_R$ is determined by
    the profiles of $\kappa$'s discriminators on
    $A \times B$. Over $\GF(2)$, a linear map is
    determined by its values on a basis.
  \item \emph{Pullback}: the relation is recovered as
    $\chi_R^{-1}(\in_B)$ by construction.
\end{itemize}
$\checkmark$

\emph{5. At the join $\kappa_1 \vee \kappa_2$.}
The join $\kappa_\vee = \mathrm{span}(\kappa_1 \cup
\kappa_2)$ contains all discriminators from both fibers.
Any composite subobject involving discriminators from
$\kappa_1$ and $\kappa_2$ is classifiable at $\kappa_\vee$.
The power object at $\kappa_\vee$ classifies all
$\kappa_\vee$-classifiable subobjects, which includes all
composites from both fibers. $\checkmark$

\emph{6. The Galois perspective.}
The short exact sequence $0 \to \kappa \to V \to
V/\kappa \to 0$ separates classifiable from
unclassifiable subobjects. $\mathrm{Gal}(V/\kappa)
= GL(V/\kappa, \GF(2))$ acts on the unclassifiable
ones. As $\kappa$ grows (via joins), $\mathrm{Gal}$
shrinks and more subobjects become classifiable.
The transition functors between fibers' power objects
are equivariant under Galois group reduction (the same
compatibility condition as in the Yoneda pro-system,
\S\ref{app:limits-exhaustive}).

In the fibered topos $\int \mathcal{C}$, the power
object at stage $\kappa$ classifies all
$\kappa$-classifiable subobjects. The fibered
construction takes the colimit over all stages,
classifying every subobject that is classifiable at
\emph{any} stage. By the operationalist axiom
(Definition \ref{def:operationalist}), a subobject
not classifiable at any stage is not a valid subobject.
$\checkmark$


%----------------------------------------------------------------------
\subsection{Sheaf topology: non-triviality verification}
\label{app:topology-nontrivial}
%----------------------------------------------------------------------

We verify that the Grothendieck topology generated by the
six triadic adjunctions is strictly between trivial and
discrete, by exhaustive enumeration of covering status.

\emph{Setup.}
The six adjunctions are organized as three axes (the three
non-zero elements of $\GF(2)^2$) with two chiralities
each (inclusion $\iota$ and projection $\pi$). Each axis
contributes a 1-dimensional sub-regime of the
2-dimensional regime $\GF(2)^2$:
\begin{center}
\begin{tabular}{c|c|c}
Axis & Element of $\GF(2)^2$ & Sub-regime \\
\hline
$\tau$ & $(1,0)$ & $\mathrm{span}\{(1,0)\}$ \\
$\sigma$ & $(0,1)$ & $\mathrm{span}\{(0,1)\}$ \\
$\kappa$ & $(1,1)$ & $\mathrm{span}\{(1,1)\}$
\end{tabular}
\end{center}

A family of sub-regimes \emph{covers} $\GF(2)^2$ iff the
sub-regimes jointly span the full 2-dimensional space.

\emph{Key fact}: any two of $\{(1,0), (0,1), (1,1)\}$
are linearly independent over $\GF(2)$ (each pair spans
$\GF(2)^2$). This is because $(1,0) + (0,1) = (1,1)$:
the three elements form the triangle identity, and any
two determine the third.

\emph{Exhaustive classification.}

\emph{1. Single-axis families}: a family using
sub-regimes from only one axis (e.g.\ both chiralities
of $\tau$). The span is 1-dimensional. Does \textbf{not}
cover. (3 such families.) $\checkmark$

\emph{2. Two-axis families}: a family using sub-regimes
from two different axes (e.g.\ $\tau$ and $\sigma$).
The span is 2-dimensional (any two of the three elements
are linearly independent). \textbf{Covers.}
($\binom{3}{2} = 3$ axis-pairs, each with $2 \times 2 = 4$
chirality choices $= 12$ families.) $\checkmark$

\emph{3. Three-axis families}: spans 2 dimensions
(contains a two-axis subfamily). \textbf{Covers.}
($2^3 = 8$ chirality choices.) $\checkmark$

\emph{Not trivial}: the two-axis families are proper
subfamilies of the full six adjunctions, and they cover.
In the trivial topology, only the maximal sieve covers.
Here, smaller sieves cover. $\checkmark$

\emph{Not discrete}: the single-axis families do not
cover. In the discrete topology, every non-empty sieve
covers. Here, some non-empty sieves fail to cover.
$\checkmark$

\emph{Minimal covering families}: pairs of sub-regimes
from different axes. There are 12 such pairs (3 axis-pairs
$\times$ 4 chirality choices). The $S_3$ symmetry
($GL(2, \GF(2)) \cong S_3$, the Galois group) acts
transitively on the three axis-pairs. The covering
condition is $S_3$-invariant: if a family covers, so
does its image under any $S_3$ permutation.

The topology is the \emph{unique} $S_3$-invariant
topology on the site of $\GF(2)^2$ sub-regimes that
declares ``two different axes cover.'' It is determined
by the triangle identity (any two vertices determine the
third) and the Galois group ($S_3$ acts transitively on
axis-pairs). $\checkmark$


%----------------------------------------------------------------------
\subsection{$(n,1)$-structure: no direction at grade $\geq 2$}
\label{app:no-direction-grade2}
%----------------------------------------------------------------------

We prove that no grade-$k$ element ($k \geq 2$) can
induce a non-invertible natural transformation. Direction
is confined to grade 1.

\emph{What direction requires.}
A non-invertible morphism $f: A \to B$ requires
\emph{asymmetry}: the forward morphism $f$ is not the
same as the reverse $f^{-1}$. At grade 1, this asymmetry
comes from $\tau/\sigma$ labels on the data: $d_f$
witnesses $A \to B$ (with $\tau_{d_f}(A) = 1$,
$\sigma_{d_f}(B) = 1$), while the reverse $B \to A$
would require a different discriminator $d'$ with
$\tau_{d'}(B) = 1$, $\sigma_{d'}(A) = 1$. In general
$d_f \neq d'$, so $f$ is non-invertible. $\checkmark$

\emph{Why grade $\geq 2$ is symmetric.}
A grade-$k$ element is $\alpha = d_1 \wedge d_2 \wedge
\cdots \wedge d_k$. Over $\GF(2)$, the wedge product is
commutative: any permutation $\pi \in S_k$ gives
\[
  d_{\pi(1)} \wedge \cdots \wedge d_{\pi(k)}
  = d_1 \wedge \cdots \wedge d_k = \alpha.
\]
The ``reverse'' of $\alpha$ --- obtained by reversing the
factor order --- is $d_k \wedge \cdots \wedge d_1 = \alpha$.
Every grade-$k$ element is its own reverse. Therefore
$\alpha^{-1} = \alpha$: every $k$-morphism ($k \geq 2$)
is self-inverse, hence invertible. $\checkmark$

\emph{Why grade 1 escapes.}
A grade-1 element $d$ has no factors to permute (only
one). Its asymmetry comes not from factor ordering but
from the $\tau/\sigma$ labels on the \emph{population}:
$d$ directly evaluates on elements of $S$, and the
evaluation $d(a)$ determines whether $a$ carries the
$\tau$ or $\sigma$ property. This is the unique grade
where the discriminator contacts the data.

At grade $\geq 2$, the element does not contact the
population directly. It is a \emph{comparison} of
discriminators, and the comparison is symmetric because
$\wedge$ is commutative. The $\tau/\sigma$ labels on
individual factors are present but their ordering is
erased by commutativity.

\emph{The Galois perspective.}
The permutation group $S_k$ acts on the $k$ factors of
a grade-$k$ element. Over $\GF(2)$, $S_k$ acts
\emph{trivially} (all signs are $+1$). The ``Galois
group of factor ordering'' is the full $S_k$, acting
trivially: there is no information about which factor
came first. Direction requires breaking this symmetry.
At grade 1, the $\tau/\sigma$ population interface
breaks it (the discriminator contacts asymmetrically
labeled data). At grade $\geq 2$, nothing breaks it.

This is why $\mathcal{C}(S,\kappa)$ is an
$(n,1)$-category: directed (non-invertible) morphisms
exist only at grade 1. All higher morphisms are
invertible, forming a groupoid at each level $\geq 2$.
$\checkmark$


%----------------------------------------------------------------------
\subsection{Dynamics transfer: linearity under $S_3$ rotation}
\label{app:dynamics-s3}
%----------------------------------------------------------------------

We verify that profile linearity and evaluation linearity
are preserved under $S_3$ rotation of perspectives.

\emph{The evaluation map.}
Profile linearity (Definition \ref{def:profile-linear})
and evaluation linearity (Definition \ref{def:eval-linear})
jointly assert that the evaluation map
\[
  \mathrm{ev}: V \times S \to \GF(2), \qquad
  \mathrm{ev}(d, a) = d(a)
\]
is \emph{bilinear}: linear in $d$ (for fixed $a$) and
linear in $a$ (for fixed $d$). When $S \cap V \neq
\varnothing$, this is the standard pairing
$\langle d, a \rangle$ on $V$.

\emph{What $S_3$ does.}
Under the $\sigma$-perspective (Proposition
\ref{prop:perspectives}), the roles rotate: what was the
discriminator space ($V$, playing $\kappa$) becomes the
population (playing $\sigma$), and what was the population
($S$, playing $\tau$) becomes the discriminator space
(playing $\kappa$). The evaluation map at the rotated
level is:
\[
  \mathrm{ev}'(a, d) = \mathrm{ev}(d, a) = d(a).
\]
This is $\mathrm{ev}$ with its arguments swapped.

\emph{Bilinearity is preserved under argument swap.}
If $\mathrm{ev}(d, a)$ is linear in $d$ (profile
linearity) and linear in $a$ (evaluation linearity),
then $\mathrm{ev}'(a, d) = \mathrm{ev}(d, a)$ is
linear in $a$ (now the first argument --- this is the
\emph{original} evaluation linearity) and linear in $d$
(now the second argument --- this is the \emph{original}
profile linearity).

The two linearities \emph{swap} under rotation:
\begin{center}
\begin{tabular}{l|ll}
& First argument & Second argument \\
\hline
Original: $\mathrm{ev}(d, a)$ &
  Profile linearity &
  Evaluation linearity \\
Rotated: $\mathrm{ev}'(a, d)$ &
  Evaluation linearity &
  Profile linearity
\end{tabular}
\end{center}

Since both hold at the original level, both hold at the
rotated level. The swap is the $S_3$ action on the axiom
pair. $\checkmark$

\emph{Full $S_3$.}
The six elements of $S_3$ permute $\{\kappa, \sigma,
\tau\}$. Each permutation relabels which component plays
which role. The evaluation map $\mathrm{ev}(d, a) =
\langle d, a \rangle$ is the standard bilinear pairing.
Relabeling does not change bilinearity: it may swap the
two linearity conditions, but both are maintained at
every perspective. $\checkmark$

\emph{The Galois perspective.}
The bilinear pairing $\langle \cdot, \cdot \rangle:
V \times V^* \to \GF(2)$ is invariant under $GL(V)$.
$S_3$ acts as a subgroup of $GL(V)$ (permuting basis
elements). Bilinearity is preserved under any linear
change of basis, hence under any $S_3$ element.
Profile linearity and evaluation linearity are two
facets of one $GL$-invariant structure: the bilinear
pairing. $S_3$ rotates the facets but preserves the
structure. $\checkmark$


%----------------------------------------------------------------------
\subsection{Finite choice: explicit termination bound}
\label{app:choice-termination}
%----------------------------------------------------------------------

We provide the explicit decreasing measure for the finite
choice argument (Theorem \ref{thm:choice}(i)).

\emph{The measure.}
$m(\kappa) = \dim(V) - \dim(\kappa) = \dim(V/\kappa)$:
the dimension of the complement. Equivalently, $m$ is
the dimension of the Galois group's domain (the space
of unseen discriminators).

\emph{Properties.}
\begin{enumerate}[label=(\roman*),nosep]
  \item $m(\kappa) \in \mathbb{N}$
    (non-negative integer).
  \item Each $\kappa$-evolution adds one independent
    discriminator (the $e_{n+1}$ of the filling procedure,
    \S\ref{app:exhaustivity-full}). Therefore
    $\dim(\kappa') = \dim(\kappa) + 1$ and
    $m(\kappa') = m(\kappa) - 1$. The measure
    \emph{strictly decreases} by exactly 1 per evolution.
  \item $m(\kappa) = 0$ iff $\kappa = V$ (every
    discriminator is present). At $\kappa = V$, every
    element of $S$ has a unique profile in $\GF(2)^n$.
    Equivalence classes are singletons; selection from
    singletons is trivial.
  \item The process terminates in at most $n = \dim(V)$
    evolutions.
\end{enumerate}

\emph{The selection procedure.}
Given a finite family $\{S_1, \ldots, S_m\}$:
\begin{enumerate}[label=\arabic*.,nosep]
  \item For each $S_i$, attempt selection. If the
    $\kappa$-equivalence classes of $S_i$ are singletons,
    selection is trivial (each class contains exactly one
    representative). Done for this $S_i$.
  \item If selection fails (some class contains multiple
    unresolved elements), the failure produces a residue
    $R_i$ (elements with $\chi = 0$).
  \item The residue $R_i$ is a $k$-cycle. The filling
    procedure (\S\ref{app:exhaustivity-full}) adds one
    independent discriminator $e_{n+1}$, producing
    $\kappa' = \kappa \vee \{e_{n+1}\}$.
  \item The new discriminator refines \emph{all}
    sub-families simultaneously (the evolution is global,
    not per-$S_i$). $m(\kappa')= m(\kappa) - 1$.
  \item Repeat from step 1 with $\kappa'$.
\end{enumerate}

\emph{Termination.}
After at most $n$ iterations, $m(\kappa) = 0$ and
$\kappa = V$. All equivalence classes are fully refined.
Selection succeeds for every $S_i$.

\emph{Bound.} The total number of $\kappa$-evolutions
for a finite family of any size is at most $n = \dim(V)$,
because evolutions are shared across sub-families and
each evolution decreases $m$ by exactly 1.

\emph{Galois perspective.}
$m(\kappa)$ is the log (base 2) of the lower bound on
$|\mathrm{Gal}(V/\kappa)|$. Each evolution reduces the
Galois group. At $m = 0$, the Galois group is trivial
--- no symmetry of the unknown remains, and every
element is uniquely determined. The selection function
is the trivial map from each singleton class to its
unique element. $\checkmark$


%----------------------------------------------------------------------
\subsection{Pairing: equivalence classes, not characteristic
  functions}
\label{app:pairing-linearity}
%----------------------------------------------------------------------

The reviewer's concern: the characteristic function
$\chi_{\{a,b\}}(x) = 1$ iff $x \in \{a,b\}$ is not
a linear functional for $|S| > 2$. Over $\GF(2)^n$
with $n \geq 2$, a linear functional $f: V \to \GF(2)$
has $|f^{-1}(1)| = 2^{n-1}$, which is $\geq 2$. A
linear functional cannot single out exactly 2 elements
from a population of $> 4$.

\emph{Resolution}: pairing does not require a linear
characteristic function. It uses $\kappa$-equivalence
classes.

\emph{The pair as an equivalence class.}
The pair $\{a, b\}$ is the $\kappa$-equivalence class
$[a]_\kappa = [b]_\kappa$ under a regime $\kappa$ where
$a$ and $b$ are indistinguishable but both are
distinguished from all other elements.

The condition: for all $d \in \kappa$, $d(a) = d(b)$.
Over $\GF(2)$, this means $a + b \in \kappa^\perp$
(the annihilator of $\kappa$ --- the subspace of vectors
killed by every functional in $\kappa$). Simultaneously,
for each $c \notin \{a, b\}$, there exists $d \in \kappa$
with $d(a) \neq d(c)$, i.e., $a + c \notin \kappa^\perp$.

\emph{Existence.} Given $a \neq b$, choose $\kappa$ such
that $a + b \in \kappa^\perp$ and $a + c \notin
\kappa^\perp$ for each $c \notin \{a, b\}$. This is a
condition on $\kappa$: include discriminators that
separate $a$ from each $c \neq b$, but exclude
discriminators that separate $a$ from $b$.

Concretely: let $\kappa = \{d \in V^* : d(a + b) = 0\}
\cap \{d : d(a + c) = 1 \text{ for some } c \notin
\{a,b\}\}$. The first condition ensures $a \sim_\kappa b$.
The second ensures $a \not\sim_\kappa c$ for $c \neq b$.
Such a $\kappa$ exists because $a + b \neq a + c$ for
$c \neq b$ (since $b \neq c$), so the conditions are
compatible. $\checkmark$

\emph{Concrete model verification.}
In the $\GF(2)^3$ model (\S\ref{sec:concrete-model}):
under $\kappa = \{e_1\}$, elements $a$ and $b$ both have
$e_1(a) = e_1(b) = 1$. If no other element has $e_1 = 1$,
then $\{a, b\} = [a]_{\{e_1\}} = \{x : e_1(x) = 1\}$.
The pair is the equivalence class. $\checkmark$

\emph{Galois perspective.}
The pair $\{a, b\}$ exists when $a + b$ is in the
annihilator $\kappa^\perp$. The annihilator is the
``Galois-fixed'' subspace: elements invisible to $\kappa$.
The pair is defined by what $\kappa$ \emph{cannot}
distinguish, not by what a predicate positively
identifies. Pairing is a $\kappa$-deformation that
deliberately withholds a discriminator, preserving the
equivalence. $\checkmark$


%----------------------------------------------------------------------
\subsection{Fixed-point stability and uniqueness}
\label{app:fixed-point-stability}
%----------------------------------------------------------------------

We prove the fixed point (Theorem \ref{thm:entailed})
is stable by showing it is the \emph{unique} non-zero
grade-2 element, and that perturbation to zero triggers
reinstatement.

\emph{Dimensionality.}
The toroid is $\GF(2)^2$, with $\dim = 2$. The grade-2
part of the exterior algebra is
$\bigwedge^2(\GF(2)^2)$, which has dimension
$\binom{2}{2} = 1$. There is exactly \textbf{one}
non-zero grade-2 element: $e_1 \wedge e_2$ (the top
form, where $e_1, e_2$ is any basis of $\GF(2)^2$).

\emph{Uniqueness.}
The 2-cell filling the toroid's triangle is a non-zero
element of $\bigwedge^2(\GF(2)^2)$. Since this space is
1-dimensional, the filling is unique (up to the scalar
$\GF(2)^* = \{1\}$, which is trivial). There is no
other non-zero 2-cell to perturb to.

This resolves the uniqueness question [O5]: the
conjecture that $H_2 = 0$ is not needed. Uniqueness
follows from $\dim \bigwedge^2(\GF(2)^2) = 1$.
$\checkmark$

\emph{Stability.}
Two possible perturbations:
\begin{enumerate}[label=(\roman*),nosep]
  \item Perturbation to another non-zero 2-cell:
    \textbf{impossible}. The space is 1-dimensional;
    there is no other non-zero element.
  \item Perturbation to zero (removing the 2-cell):
    triggers the four failures of Theorem
    \ref{thm:irremovable} (sheaf, manifold, bimodule,
    self-resolution). The homological gap produced by
    removal specifies its own filling, which is the
    unique non-zero 2-cell $e_1 \wedge e_2$. Removal
    reinstates.
\end{enumerate}
The fixed point is stable: any departure returns to
$e_1 \wedge e_2$. $\checkmark$

\emph{$GL$-invariance.}
$GL(2, \GF(2)) \cong S_3$ acts on $\GF(2)^2$. The top
form $e_1 \wedge e_2$ transforms under $g \in GL(2)$ as
$g(e_1) \wedge g(e_2) = \det(g) \cdot (e_1 \wedge e_2)$.
Over $\GF(2)$, $\det: GL(2, \GF(2)) \to \GF(2)^* = \{1\}$
(the determinant is always 1, since $\GF(2)^*$ has one
element). Therefore $e_1 \wedge e_2$ is \emph{invariant}
under the full Galois group $S_3$. The 2-cell is the
unique $S_3$-invariant grade-2 element. $\checkmark$

\emph{Summary.}
The fixed point is:
\begin{itemize}[nosep]
  \item \textbf{unique} (1-dimensional top form),
  \item \textbf{stable} (perturbation to zero reinstates;
    perturbation to another element is impossible),
  \item \textbf{$S_3$-invariant} (determinant is trivial
    over $\GF(2)$).
\end{itemize}
These three properties follow from $\dim \bigwedge^2
(\GF(2)^2) = 1$ and $|\GF(2)^*| = 1$. $\checkmark$


%----------------------------------------------------------------------
\subsection{Foundation chain bound: link independence}
\label{app:chain-bound}
%----------------------------------------------------------------------

We prove that distinct membership links in an acyclic
chain require linearly independent witnessing
discriminators, bounding chain depth by $\dim(\kappa)$.

\emph{Link vectors.}
In a chain $a_1 \ni a_2 \ni \cdots \ni a_m$
(equivalently $a_2 \in a_1,\; a_3 \in a_2, \ldots$),
each link $a_i \ni a_{i+1}$ is witnessed by a
discriminator $d_i$. Define the \emph{link vector}
$v_i = a_i + a_{i+1} \in \GF(2)^n$ (the profile
difference; over $\GF(2)$, difference $=$ sum).

Each $v_i$ is non-zero (since $a_i \neq a_{i+1}$ in an
acyclic chain). Each $v_i$ is detected by $d_i$:
$d_i(v_i) = d_i(a_i) + d_i(a_{i+1}) = 1$ (since $d_i$
distinguishes $a_i$ from $a_{i+1}$). $\checkmark$

\emph{Independence requirement.}
If $v_i$ is linearly dependent on previous link vectors
$\{v_1, \ldots, v_{i-1}\}$, then
$v_i = \sum_{j < i} c_j v_j$ for some
$c_j \in \GF(2)$. This means:
\[
  a_i + a_{i+1} = \sum_{j < i} c_j (a_j + a_{j+1}).
\]
The profile difference at link $i$ is a $\GF(2)$-linear
combination of previous profile differences. The
distinction between $a_i$ and $a_{i+1}$ is therefore
\emph{derivable} from the distinctions already
established by earlier links. No new discriminator
is needed: the existing discriminators in $\kappa$
already detect $v_i$ (since they detect each $v_j$
and $v_i$ is in their span).

A link whose vector is dependent on previous links does
not deepen the chain --- it is a derived relation, not
a primitive membership step. The \emph{depth} of the
chain (the number of genuinely new distinctions) equals
$\dim(\operatorname{span}\{v_1, \ldots, v_{m-1}\})$.

\emph{The bound.}
The link vectors $v_1, \ldots, v_{m-1}$ lie in
$\GF(2)^n$, which has dimension $n$. Therefore:
\[
  \text{chain depth}
  = \dim(\operatorname{span}\{v_1, \ldots, v_{m-1}\})
  \leq n = \dim(V).
\]
An acyclic chain has at most $n$ independent links
and at most $n + 1$ elements.

\emph{Galois measure.}
Each independent link vector $v_i$ requires a
discriminator $d_i$ with $d_i(v_i) = 1$. Independent
link vectors require linearly independent discriminators
(if $d_i = d_j$ then $d_i$ detects both $v_i$ and $v_j$
only if $v_i$ and $v_j$ lie outside $\ker(d_i)$, but
this does not constrain $v_i$ to be independent of
$v_j$; however, if $v_i \in \operatorname{span}\{v_j\}$
then link $i$ is derived, as above).

The measure $m(\kappa) = \dim(V/\kappa)$ decreases with
each genuinely independent link:
adding $d_i$ to $\kappa$ for an independent $v_i$
increases $\dim(\kappa)$ by 1 and decreases $m$ by 1.
After $n$ independent links, $m = 0$, the Galois group
is trivial, and no further independent links are
possible.

\emph{Sharp bound.}
The bound $n + 1$ is achieved: in $\GF(2)^n$ with
basis $e_1, \ldots, e_n$, the chain
$0 \ni e_1 \ni (e_1 + e_2) \ni \cdots \ni
(e_1 + \cdots + e_n)$ has $n$ links with link vectors
$e_1, e_2, \ldots, e_n$ (a basis). $\checkmark$


%----------------------------------------------------------------------
\subsection{Indistinguishability: $\omega$-chain stabilization}
\label{app:omega-chain}
%----------------------------------------------------------------------

We rule out convergent $\omega$-chains of
$\kappa$-evolutions separating actual from potential
infinity.

\emph{The concern.}
Proposition \ref{prop:indistinguish} shows each
individual discriminator fails to separate ``actually
infinite'' from ``potentially infinite (finite but
unbounded).'' The reviewer asks: could an $\omega$-chain
$\kappa_1 \subseteq \kappa_2 \subseteq \cdots$ produce
a \emph{collective} separation, even though no individual
discriminator separates?

\emph{Step 1: stabilization.}
$V$ is finite-dimensional ($\dim V = n$). The ascending
chain $\kappa_1 \subseteq \kappa_2 \subseteq \cdots$
is a chain of subspaces of $V$. By the ascending chain
condition in a finite-dimensional vector space, the chain
stabilizes: there exists $N$ such that $\kappa_N =
\kappa_{N+1} = \cdots$. The ``limit'' regime is
$\kappa_\infty = \kappa_N$, which has
$\dim(\kappa_\infty) \leq n$. $\checkmark$

\emph{Step 2: finite span.}
Every element of $\kappa_\infty$ is a linear combination
of finitely many $d_i$'s (since $\kappa_\infty$ is
finite-dimensional and the $d_i$ span it). No ``limit
discriminator'' exists outside $\operatorname{span}
\{d_i\}$: over a finite-dimensional space, the span
is all there is. $\checkmark$

\emph{Step 3: linearity kills collective separation.}
The original argument shows: each $d_i$ individually
gives the same result on ``actually infinite'' and
``potentially infinite'' sequences. That is,
$d_i(\text{actual}) = d_i(\text{potential})$ for all $i$.

By profile linearity, any linear combination
$d = \sum c_i d_i$ ($c_i \in \GF(2)$) satisfies:
\[
  d(\text{actual}) = \sum c_i\, d_i(\text{actual})
  = \sum c_i\, d_i(\text{potential})
  = d(\text{potential}).
\]
A linear combination of non-separating discriminators
is non-separating. Therefore no element of
$\kappa_\infty$ separates. $\checkmark$

\emph{Conclusion.}
No $\omega$-chain of $\kappa$-evolutions converges to a
separation of actual from potential infinity. The
stabilization (step 1) ensures the chain has a finite
limit. The finite-dimensionality (step 2) ensures no
discriminator lies outside the span. The linearity
(step 3) ensures nothing in the span separates.

\emph{Galois perspective.}
The descending chain $\mathrm{Gal}(V/\kappa_1) \supseteq
\mathrm{Gal}(V/\kappa_2) \supseteq \cdots$ is a chain
of subgroups of the finite group $GL(V, \GF(2))$.
It stabilizes at $\mathrm{Gal}(V/\kappa_\infty)$.
The residual Galois group contains the symmetries that
\emph{cannot} be broken by any $\kappa$-evolution ---
and the actual/potential distinction lies within these
unbreakable symmetries. No amount of Galois group
reduction reaches it, because it is not a distinction
that any $\GF(2)$-linear discriminator can witness.
$\checkmark$


%----------------------------------------------------------------------
\subsection{Native univalence from the annihilator
  intersection}
\label{app:native-univalence}
%----------------------------------------------------------------------

We construct the precise univalence statement for the
discrimination framework.

\emph{The path space.}
For two elements $a, b \in S$, the ``path'' from $a$ to
$b$ is the vector $a + b \in V$ (over $\GF(2)$,
difference $=$ sum). The elements $a$ and $b$ are
$\kappa$-equivalent iff $a + b \in \kappa^\perp$ (the
annihilator of $\kappa$: the subspace of vectors killed
by every discriminator in $\kappa$).

The \emph{path space} between $[a]$ and $[b]$ across all
stages is:
\[
  \mathrm{Path}(a, b) = \{v \in V : v = a + b,\;
  v \in \kappa^\perp \text{ for all } \kappa
  \text{ where } a \sim_\kappa b\}.
\]
For persistent equivalence (equivalence at every stage):
$\mathrm{Path}(a, b) = \{a + b\} \cap
\bigcap_\kappa \kappa^\perp$.

\emph{The operationalist axiom as $\bigcup \kappa = V$.}
The operationalist axiom (Definition
\ref{def:operationalist}) states: if a distinction cannot
be witnessed by any discriminator in any $\kappa$-state,
it is not a valid distinction. Algebraically: the union
of all $\kappa$-regimes spans $V$:
\[
  \bigcup_\kappa \kappa = V.
\]
Every direction in $V$ is eventually probed by some
discriminator.

\emph{Annihilator intersection.}
Taking annihilators reverses inclusions. If
$\bigcup_\kappa \kappa = V$, then:
\[
  \bigcap_\kappa \kappa^\perp
  = \left(\bigcup_\kappa \kappa\right)^\perp
  = V^\perp = \{0\}.
\]
The only vector invisible to \emph{all} stages is the
zero vector. $\checkmark$

\emph{Univalence.}
If $a \sim_\kappa b$ for all $\kappa$ (persistent
equivalence), then $a + b \in \bigcap_\kappa
\kappa^\perp = \{0\}$. Therefore $a + b = 0$, i.e.,
$a = b$.

The \textbf{ua map}: persistent equivalence $\to$
identity is $a + b \mapsto 0$ (the equivalence data
forces $a + b = 0$, which is $a = b$).

The \textbf{idtoeqv map}: identity $\to$ equivalence
is $a = b \implies a + b = 0 \in \kappa^\perp$
for all $\kappa$ (the zero vector is in every
annihilator).

These are inverses:
$\mathrm{ua} \circ \mathrm{idtoeqv} = \mathrm{id}$
(identity implies equivalence implies $a + b = 0$
implies identity);
$\mathrm{idtoeqv} \circ \mathrm{ua} = \mathrm{id}$
(equivalence implies $a + b = 0$ implies identity
implies equivalence).
$\checkmark$

\emph{Contractibility of the path space.}
For $a = b$: $\mathrm{Path}(a, a) = \{0\}$, a single
point. Contractible.
For $a \neq b$ but $a \sim_\kappa b$ at some stage:
$\mathrm{Path}(a, b) = \{a + b\}$ with $a + b \neq 0$.
This is a single non-zero point at stage $\kappa$, but
as $\kappa$ grows, eventually some discriminator detects
$a + b$ (since $\bigcup \kappa = V$), and the equivalence
breaks. The path space empties.

In both cases, the non-empty path space is contractible
(a single point). This is the HoTT condition for
univalence: the fiber of $\mathrm{idtoeqv}$ over each
equivalence is contractible. $\checkmark$

\emph{The Galois reading.}
At each stage, $\kappa^\perp$ is the Galois group's
domain: the directions invisible to $\kappa$. The path
$a + b$ lives in $\kappa^\perp$ iff $a$ and $b$ are in
the same Galois orbit. As $\kappa$ grows, the Galois
group shrinks and orbits refine. Persistent equivalence
means $a + b$ is in every Galois group's domain. By
the operationalist axiom, the intersection of all
domains is $\{0\}$. The only persistent orbit is the
trivial orbit. Univalence is Galois-group exhaustion.
$\checkmark$


%----------------------------------------------------------------------
\subsection{Segal conditions at $k \geq 3$: groupoid
  filling}
\label{app:segal}
%----------------------------------------------------------------------

We verify the Segal conditions for the $(n,k)$-structure
at $k \geq 3$.

\emph{The Segal condition.}
A simplicial object $X_\bullet$ satisfies the Segal
condition at level $k$ if the Segal map
\[
  X_k \to X_1 \times_{X_0} X_1 \times_{X_0}
  \cdots \times_{X_0} X_1 \quad (k \text{ factors})
\]
is an equivalence: a $k$-cell is determined by its
$k$ composable edges.

\emph{Fact 1: groupoid at grade $\geq 2$}
(Appendix \ref{app:no-direction-grade2}).
Every grade-$k$ element ($k \geq 2$) is self-inverse:
$d_{\pi(1)} \wedge \cdots \wedge d_{\pi(k)} =
d_1 \wedge \cdots \wedge d_k$ for all $\pi \in S_k$
(commutativity of $\wedge$ over $\GF(2)$). All
$k$-morphisms are invertible. $\checkmark$

\emph{Fact 2: groupoids satisfy Segal conditions.}
In a groupoid, every $k$-cell is uniquely determined
by its boundary of composable 1-cells:
\begin{itemize}[nosep]
  \item \emph{Existence}: given $k$ composable 1-forms
    $d_1, \ldots, d_k$ (pairwise linearly independent),
    their wedge product $d_1 \wedge \cdots \wedge d_k$
    is a non-zero $k$-cell.
  \item \emph{Uniqueness}: the wedge product of a given
    set of factors is unique (one product per set, since
    $\wedge$ is commutative and $\GF(2)^* = \{1\}$).
  \item \emph{No filling ambiguity}: in a groupoid,
    composition and inverses are unique. There is no
    choice in how to fill a $k$-simplex given its edges.
\end{itemize}
The Segal map is a bijection. $\checkmark$

\emph{The Grassmannian structure.}
Decomposable grade-$k$ elements biject with
$k$-dimensional subspaces of $V$ (the Grassmannian
$\mathrm{Gr}(k, V)$ over $\GF(2)$). A $k$-subspace
is determined by any basis, and all bases yield the
same wedge product (since $\GF(2)^* = \{1\}$ makes
the determinant trivial).

The Segal map at level $k$:
$\mathrm{Gr}(k, V) \to \{\text{$k$-tuples of independent
1-forms}\}/\sim$
is a bijection: a $k$-subspace $\leftrightarrow$ an
unordered set of $k$ independent generators.

The number of $k$-cells is the Gaussian binomial
coefficient $\binom{n}{k}_2$, which counts
$k$-subspaces of $\GF(2)^n$. $\checkmark$

\emph{Explicit verification at $k = 3$.}
$d_1 \wedge d_2 \wedge d_3$ is determined by
$\{d_1, d_2, d_3\}$ (unordered). Its three 2-faces
are $d_1 \wedge d_2$, $d_2 \wedge d_3$,
$d_1 \wedge d_3$. Given these three 2-faces, the
3-cell is their unique filler (exists by groupoid
composition, unique by commutativity of $\wedge$).

At $k = 4$: six 2-faces, four 3-faces, one 4-cell.
The 4-cell $d_1 \wedge d_2 \wedge d_3 \wedge d_4$ is
determined by any of its 3-faces (which are themselves
determined by their 2-faces). The Segal condition
holds inductively. $\checkmark$

\emph{Why this was non-obvious.}
The concern was that non-decomposable elements (sums
of wedge products, e.g.\ $d_1 \wedge d_2 + d_3 \wedge
d_4$) might violate Segal. But non-decomposable
elements are \emph{chain-group elements} (formal sums
of $k$-morphisms), not individual $k$-morphisms. The
Segal condition applies to the $k$-morphisms themselves
(decomposable elements = points of the Grassmannian),
not to their formal sums. The exterior algebra's full
grade-$k$ component $\bigwedge^k V$ is the chain group;
the Segal space is the Grassmannian
$\mathrm{Gr}(k, V) \subset \bigwedge^k V$. $\checkmark$


%----------------------------------------------------------------------
\subsection{O1 and O4: the RISC instruction set}
\label{app:risc}
%----------------------------------------------------------------------

O1 (scope of evaluation linearity) and O4
(join-semilattice of $\kappa$-states) are the same
question, and both close simultaneously.

\emph{O1: characterization of admissible discriminators.}
Profile linearity (Definition \ref{def:profile-linear})
requires $d(a + b) = d(a) + d(b)$. Evaluation linearity
(Definition \ref{def:eval-linear}) requires the same in
the second argument. Together: $\mathrm{ev}(d, a) =
\langle d, a \rangle$ is bilinear.

A function $f: \GF(2)^n \to \GF(2)$ satisfying
$f(a + b) = f(a) + f(b)$ is a $\GF(2)$-linear
functional. These are exactly the elements of the dual
space $V^* = \operatorname{Hom}(V, \GF(2))$.

The admissible discriminators are $V^*$. $\checkmark$

\emph{Count.} $\dim(V^*) = \dim(V) = n$. $|V^*| = 2^n$
(including zero). Non-trivial: $2^n - 1$. At $n = 3$:
$8$ admissible (including zero), $7$ non-trivial. These
$7$ are the points of $PG(2, \GF(2))$, the Fano plane.

The ratio $2^n / 2^{2^n}$ (admissible over total) is the
compression ratio: $3.1\%$ at $n = 3$, $0.024\%$ at
$n = 4$. The ``RISC instruction set'' is the projective
space $PG(n{-}1, \GF(2))$: a minimal, uniform set of
primitives from which $\kappa$-evolution recovers full
expressiveness.

\emph{O4: the join-semilattice.}
$\kappa$-states are subspaces of $V^*$. $V^*$ is a
vector space. The join of two subspaces is their span:
$\kappa_1 \vee \kappa_2 = \operatorname{span}(\kappa_1
\cup \kappa_2)$. The span of subspaces of a vector
space is a subspace.

Therefore:
\begin{itemize}[nosep]
  \item The join is well-defined and produces a valid
    $\kappa$-state. $\checkmark$
  \item Profile linearity is preserved: every element of
    $\operatorname{span}(\kappa_1 \cup \kappa_2)$ is a
    $\GF(2)$-linear combination of linear functionals,
    hence a linear functional. $\checkmark$
  \item The meet $\kappa_1 \wedge \kappa_2 = \kappa_1
    \cap \kappa_2$ is also a subspace. $\checkmark$
\end{itemize}

The lattice of $\kappa$-states is not merely a
join-semilattice. It is a \emph{modular lattice}
(the lattice of subspaces of a finite-dimensional
vector space), which is stronger than what O4 requires.
$\checkmark$

\emph{Why they are the same question.}
O1 asks: what are the primitives? Answer: $V^*$ (linear
functionals). O4 asks: do the primitives compose well?
Answer: $V^*$ is a vector space, so its subspace lattice
is a lattice. The admissibility of the instruction set
(O1) and the composability of regimes (O4) are both
consequences of the single fact that $V^*$ is a vector
space over $\GF(2)$.

\emph{Consequence for the topos.}
O4 was the sole remaining component of Theorem
\ref{thm:topos} (fibered topos). With O4 resolved,
all components are verified:
\begin{itemize}[nosep]
  \item Finite limits (G8, Appendix \ref{app:limits-exhaustive}). $\checkmark$
  \item Subobject classifier (B.9). $\checkmark$
  \item Power objects (G14, Appendix
    \ref{app:power-composite}). $\checkmark$
  \item Join-semilattice of $K$ (this section). $\checkmark$
\end{itemize}
The fibered category $\int \mathcal{C}$ is an elementary
topos. $\checkmark$


%----------------------------------------------------------------------
\subsection{O6: Cayley--Dickson hexagon axioms via
  in-universe recognizer construction}
\label{app:cd-recognizer}
%----------------------------------------------------------------------

The periodic table conjecture (Conjecture \ref{conj:cd})
asks whether the hexagon axioms hold at Cayley--Dickson
step 2 over $\GF(2)^2$. We resolve this by reframing
the CD construction as an in-universe recognizer
construction.

\emph{The three levels.}
\begin{description}[nosep,leftmargin=1cm]
  \item[Level 0] Half-bits. $d: S \to \GF(2)$. A single
    discriminator output. Recognizer: the discriminator
    itself. Grade 1 of $\bigwedge V$.
  \item[Level 1] Bits (the toroid). $(\tau_d(a),
    \sigma_d(a)) \in \GF(2)^2$. A membership profile.
    Coproduct of two half-bits. Recognizer: the composed
    discriminator $d_\tau \wedge d_\sigma$. Grade 2.
  \item[Level 2] Pairs of bits.
    $((\tau_1, \sigma_1), (\tau_2, \sigma_2)) \in
    \GF(2)^4 = (\GF(2)^2)^2$. Coproduct of coproducts
    of half-bits. Recognizer: a discriminator operating
    on the space of discriminators.
\end{description}

\emph{The CD framing (Option A).}
Treating $(\GF(2)^2)^2$ as a CD doubling gives a
4-dimensional algebra with multiplication
$(a,b)(c,d) = (ac + db,\; da + bc)$ (conjugation is
trivial since $-1 = 1$ over $\GF(2)$). This algebra
has zero divisors. The hexagon axioms for the resulting
braided structure are unclear.

\emph{The recognizer framing (Option B).}
A recognizer for a pair-of-bits is a discriminator
operating on the \emph{space of discriminators}. Under
the $\sigma$-rotation (Proposition \ref{prop:perspectives}),
discriminators become data and data become discriminators.
A level-2 recognizer is a grade-2 element of $\bigwedge V$
applied to grade-1 elements (discriminators-as-data).

Grade-2 elements are self-inverse (Appendix
\ref{app:no-direction-grade2}), form a groupoid
(Theorem \ref{thm:groupoid}), and their composition
is strict (Appendix \ref{app:monoidal-coherence}). The
recognizer-for-recognizers uses the same $\wedge$
product, the same bilinear pairing, the same Galois
structure as the level-1 recognizer.

\emph{Resolution.}
The hexagon axioms at step 2 hold for the same reason
they hold at step 1: strictness. The $\sigma$-rotation
maps the level-2 recognizer construction to the same
strict monoidal category whose coherences are all
identities (Appendix \ref{app:monoidal-coherence}).
The hexagon diagrams commute because every morphism in
them is the identity. $\checkmark$

\emph{The zero divisors are residues.}
The CD zero divisors $(1,0) \cdot (0,1) = (0,0)$ in
$\GF(2)^2$ are $\tau \cdot \sigma = 0$: the product
of a $\tau$-only profile with a $\sigma$-only profile
is the zero profile. This is the gap/overlap residue
$\chi = 0$, which triggers $\kappa$-evolution. What
the CD construction sees as algebraic degeneration
(zero divisors preventing division), the discrimination
framework sees as residue detection (the signal to
evolve $\kappa$). The ``obstruction'' is the feature.

\emph{Proof rotation.}
The CD framing (algebraic arm) makes the hexagon axioms
look like a new question at each doubling step. The
recognizer framing (coherence arm via $\sigma$-rotation
through the Galois hub) reveals it as the same
question at every level. This is a proof rotation
(\S\ref{sec:rotation}): the CD construction is the
algebraic arm's view of a structure that, rotated,
is the same strict monoidal category at every grade.
$\checkmark$

\emph{The Baez--Dolan correspondence.}
The periodic table of $(n,k)$-categories is generated by
in-universe recognizer construction: at each level, the
framework builds recognizers for the previous level's
constructs, using $\sigma$-rotation. The coherence at
each level is inherited from the strictness of $\wedge$.
The periodic table is not an external classification
imposed on the framework; it is the framework's own
recognizer hierarchy, viewed from the algebraic arm as
iterated CD doubling and from the coherence arm as
iterated $\sigma$-rotation of the same strict structure.
